Foucault, Michel. \textit{As palavras e as coisas}.
Neste ponto da escrita, a lembrança do prefácio citado me levou à sua releitura, interrompendo o texto que escrevia no diário. O que diria ali é que Cecinepá permite articular diferentes epistemes, diferentes tabuleiros lógicos, semânticos, cosmológicos, epistêmicos, com os quais se realizem discursos e performances de modo geral --- proposições. A arquitetura deste mundo heterotópico é tal que permite a disposição, a visualização e a articulação de proposições distintas que possuam um referencial comum orientando-se não por alguma noção de \enquote{verdade}, mas pelas diferentes relações que cada proposição encontra com o referencial comum.